% ======================================================================
% Drop-in for the manuscript: consensus-layer sensitivity study.
% Placement: Sec. VI gains subsection (methodology); Sec. VII gains
% subsection (results). Contains NO experimental numbers — every value
% is supplied exclusively by the auto-derived table_consensus.tex and
% fig_consensus_latency.tex, which analyze_consensus_comparison.py
% refuses to generate without measured evidence from both networks.
% Also add to the contributions list:
%   \item A consensus-layer sensitivity study anchoring the identical
%   audit micro-benchmark on Ethereum Sepolia (Proof-of-Stake) and the
%   BNB Smart Chain testnet (Proof-of-Staked-Authority, delegated-staking
%   BFT family), comparing execution-gas parity, inclusion latency,
%   observed block cadence, and preserved negative-security rejection.
% ======================================================================

% ---------------- Sec. VI addition (methodology) ----------------
\subsection{Consensus-Layer Sensitivity Methodology: PoS versus
Delegated-Staking Anchoring}
\label{sec:consensus-method}

The Stage~4 Sepolia micro-benchmark establishes that the audit anchor
executes identically on a public Proof-of-Stake (PoS) network, interpreted
strictly as an EVM execution-gas consistency check. To evaluate whether the
anchoring layer is additionally sensitive to the \emph{consensus family}, the
identical \texttt{Stage3AnomalyLedger} bytecode is deployed unchanged to a
second public testnet whose consensus belongs to the delegated-staking
family: the BNB Smart Chain testnet (chain~ID~97), which operates
Proof-of-Staked-Authority (PoSA), a BFT-style protocol in which a bounded
set of validators is elected through delegated staking and is therefore
commonly grouped with Delegated Proof-of-Stake (DPoS) designs. This
terminology is used precisely: the comparison is between Ethereum's
slot-based PoS (Gasper) on Sepolia and a delegated-staking PoSA chain, not
a claim about any specific pure-DPoS ledger.

For each network, one benchmark run deploys the contract and submits $R$
sequential anchor commitments with fresh commitment/evidence digests. Per
commitment, the harness records the receipt gas, block number, chain
timestamp, effective gas price, and the wall-clock interval from
transaction submission to first confirmation (inclusion latency). Observed
inter-block cadence is derived from the chain timestamps of the anchored
blocks. The three negative-security invariants --- replay rejection,
zero-anchor rejection, and unauthorized-submitter rejection --- are then
exercised on-chain via static calls using the same procedure as the
Stage~4 script, verifying that the enforcement semantics of
$P_5$ and the role gate are preserved across consensus families. All rows
are written to an append-only JSONL evidence file whose SHA-256 digest is
embedded in the run summary, and every manuscript value is derived
deterministically from these summaries by a generation script that fails
in the absence of measured evidence, so no reported number can originate
outside a recorded run.

Three boundaries apply. First, gas is consensus-independent EVM execution
cost; identical gas across networks validates execution parity, not
consensus performance. Second, inclusion latency is measured under
prevailing public-testnet load and validator configuration and is a
first-confirmation figure, not economic finality; Ethereum finality
additionally requires attestation epochs, and PoSA fast finality follows
its own rule set, neither of which is claimed here. Third, testnet gas
prices carry no economic meaning and are recorded only for completeness.
The study therefore answers a bounded question: does the accountable-anchoring
layer of ZKTrustLLM-Agents~L4 transfer unchanged --- in execution cost and
negative-security behaviour --- between PoS and delegated-staking chains,
and what inclusion-latency envelope does each family exhibit for the same
governed audit transaction?

% ---------------- Sec. VII addition (results) ----------------
\subsection{Consensus-Layer Sensitivity Results}
\label{sec:consensus-results}

Table~\ref{tab:consensus} reports the measured comparison and
Fig.~\ref{fig:consensus-latency} visualises the inclusion-latency
envelopes. Three observations are supported directly by the evidence.
First, execution parity: the anchor gas measured on the PoSA testnet is
reported alongside the Sepolia and local Hardhat values, extending the
execution-gas consistency check across a second consensus family.
% After running: state the measured parity/delta here in one sentence,
% citing only Table~\ref{tab:consensus}. \todo{one sentence, measured}
Second, the negative-security invariants --- replay, zero-anchor, and
unauthorized-submitter rejection --- are preserved on both networks,
confirming that the trust-plane enforcement of the audit ledger does not
depend on the consensus family beneath it. Third, the inclusion-latency
envelopes differ with block cadence:
% After running: one or two sentences describing the measured median/P95
% contrast and observed inter-block medians, citing only the table and
% figure. Do not attribute causes beyond block cadence without evidence.
% \todo{one to two sentences, measured}
These results are public-testnet observations under prevailing load and
must not be read as mainnet, economic-finality, or production O-RAN
performance claims; their scientific role is to show that the
evidence-anchoring layer is portable across consensus families with
unchanged execution cost and unchanged rejection semantics, while the
latency envelope --- and hence the achievable audit-anchoring cadence ---
is a deployment-time property of the selected chain.

\input{tables/table_consensus}
\input{figures/fig_consensus_latency}
